\scanpage{39}
Nous allons introduire une autre structure de nature \unsure{transcendante}, en introduisant le \unsure{rev.|ver.} universel $\bar S_{\mathbb C}$ de $S(\mathbb C)$ en $\xi$ $\unsure{\xi|:}$ [sous-entendu : de la comp. connexe de $\xi$ dans $S(\mathbb C)$].

Soit $\unsure{\gamma|\mathscr{Y}}_{G}$ l'ensemble des hom. de $\unsure{\mathbb S|S}=\prod_{\mathbb C/\mathbb R}\mathbb G_{m\,\mathbb C}$ dans $G_{\mathbb R}$, correspondant à l'hom central \deleted{\unsure{donné}} $i:\mathbb G_m \unsure{\relbar|\longrightarrow} G$ donné par $i(\lambda)=i_1(\lambda)i_2(\lambda)$. \margin{\rule{0.35pt}{2.1\baselineskip}}On considère $\gamma$ comme variété \deleted{\unsure{an.|anal.}} holomorphe où \deleted{$\unsure{\gamma|G}$} $G(\mathbb R)$ et à fortiori $G(\unsure{\mathbb Q|Q})$ opère. Ceci \unsure{posé|fait}, on a une \uline{application}
\[
\boxed{f:\bar S_{\mathbb C}\longrightarrow\gamma_G}
\]
satisfaisant les conditions

\uline{12)}\unsure{a.|.} \ a) $f(\bar\xi\inserted{\text{le pt marqué de }\bar S_{\mathbb C}})={}\text{l'hom }\mathbb S\unsure{\relbar|\longrightarrow} G_{\mathbb R}$ défini par $i_1^{\mathbb C},i_2^{\mathbb C}$ de c).

b) $f(g\bar s)=\unsure{w_{1}|w|h_{1}|h_{\gamma}|\rho_{1}|\rho_{\gamma}}(g)f(\bar s)$ si $g\in\pi_1(S(\mathbb C),\xi),\ \bar s\in\bar S_{\mathbb C}$

c) $f$ est \uline{holomorphe}.

\unsure{14)|15)} (Moyennant Hodge) \deleted{Pour tout $\bar s\in\bar S_{\mathbb C}$, le plus grand sous-groupe algébrique de $G$ qui sur $\mathbb R$ contient $f(\bar s)(\mathbb S)\subset G_{\mathbb R}$, et qui de plus contient $H$} \inserted{Le composé \deleted{\unsure{de}} $\bar S_{\mathbb C}\longrightarrow\gamma_G\longrightarrow\gamma_{\widetilde G}$} (qui à priori \inserted{\unsure{vû|vu|n\textsuperscript{o}} : 12 b)} se factorise par \deleted{\unsure{$\bar S_{\mathbb C}$}} $S(\mathbb C)^{\unsure{\circ|0}}$) est une application \uline{constante}.

13) Si $\bar s\in\bar S_{\mathbb C}$ se trouve sur le point $t$ de $S$ [il nous \unsure{plaît|plait|suffit} bien que $P$ comme torseur à connexion \unsure{intégr.} $P_S$ soit défini sur $S$] de sorte qu'on a torseur à conn. intégrable \deleted{\illegible{1}} $(P_t,G_{k(t)})$ sur $k(t)$, \deleted{\unsure{soit}} l'hom $f(\bar s)_{\mathbb C}:\mathbb S_{\mathbb C}\longrightarrow(G_{k(t)})_{\mathbb C}\simeq G_{\mathbb C}$ [définissant une bigraduation\inserted{ donc filtration} sur le foncteur $T_{\mathbb C}$, $T$ foncteur oubli sur les $G$-modules] est \uline{compatible} avec le torseur $P_t$ sous $G_{k(t)}$ et sa $\unsure{C|\mathbb C|G}$-trivialisation $p(\bar s)$ [$p:\bar S_{\mathbb C}\longrightarrow(P_S)_{\mathbb C}$ $\unsure{\in|\epsilon}$ ici est \unsure{le|la} $S(\mathbb C)^{\circ}$-morphisme \unsure{caractérisé} par la condition que $p(\bar\xi)=\alpha$].
% archive sheet number: 38
